\section{Predestination and Predilection}

\begin{quotex}
Our will is truly free only in union with that of God and that God acts on earth only through our free will freely united with his. \flright{\textsc{Valentin Tomberg}}

Every incarnated human being is the product of two shaping forces: heredity and the creative force of self-realization of the eternal individuality. \flright{\textsc{Valentin Tomberg}}

\end{quotex}
Any discussion of the mystery of birth would be incomplete without an understanding of predestination and predilection as understood in the Medieval Germanic-Roman religion, which is also known as the Catholic religion. To summarize the earlier post:

\begin{itemize}
\item Each person comes into the world with a set of qualities (his karmic inheritance). 
\item Beyond those innate qualities, his life's path is to actualize the possibilities open to him (his dharma path). 
\item He enters the world process at a time, place, and relationships suitable for him. 
\end{itemize}
Nevertheless, the story is incomplete. Ultimately, beyond such terrestrial concerns, the goal is salvation and liberation. Hence, there is a supernatural destiny beyond our ordinary life. For some reason, the doctrines of predestination and predilection are seldom mentioned today despite their enormous importance. They may be hard to understand, or even distasteful to the secular mind, so it is helpful to meditate on them. For what follows, I am relying mostly on the book \emph{Predestination} by \textbf{Fr Reginald Garrigou-Lagrange}.

\paragraph{Predilection}
This is the \textbf{Principle of Predilection}:

\emph{One would not be better than another unless one were loved more by God. }

Predilection and predestination are clearly radically anti-egalitarian, as that principle makes clear.

\textbf{St Thomas Aquinas} justifies this principle:

\begin{quotex}
Since God's love is the cause of goodness in things, no one thing would be better than another if God did not will greater good for one than for another … and the reason why some things are better than others is that God wills them a greater good. Hence it follows that He loves more the better things.

\end{quotex}
Hence good things are not better because of some accident or injustice, but precisely because God loves them more. And the most perfect good is \emph{theosis}, or becoming God-like. Fr Garrigou-Lagrange explains:

\begin{quotex}
Every agent acts for an end and the purpose of the action of the supreme agent is to manifest His goodness by reproducing a likeness of Himself which is a more or less perfect participation of His nature.

\end{quotex}
Predilection is both a philosophical truth, known to reason, as well as a revealed truth. As a philosophical truth, it is obvious that it is true. Fr. Garrigou-Lagrange explains:

\begin{quotex}
This principle is true in every order. It is true of plants, of animals, of human beings, of angels, of things in which there is less of perfection or of goodness. It is also true of every man who, from whatever point of view, is better than another …

\end{quotex}
As a revealed truth, predilection is related to the gratuitous gift of grace from God which makes us pleasing. Thus, predilection, or what we consider the “good life”, may have one meaning in the natural realm, but in the supernatural, it may be completely different; the beatitudes are examples of supernatural predilection. This goodness comes from the grace of God, and makes us pleasing in his sight.

There is no injustice in this, for justice requires that unlike things be treated differently. Nevertheless, pride must be avoided since it is God who makes one man superior to another (Thomas Aquinas).

\paragraph{Predestination}
\begin{quotex}
Predestination is the foreknowledge and preparedness on God's part to bestow the favors by which all those are saved who are to be saved … God already knew, when He predestined, what He must do to bring His elect infallibly to eternal life. \flright{\textsc{St. Augustine}}
\end{quotex}
The doctrine of predestination follows from predilection. St Thomas defines it this way:

\begin{quotex}
The plan of the direction of a rational creature towards the end, i.e., life eternal; for to destine is to direct or send.

\end{quotex}
Therefore, predestination is the plan in God's mind of directing a particular man to the ultimate and supernatural end. This plan, from all eternity, determines the efficacious means that will lead a man to his final end.

For God, intention precedes execution, or ends before means. Hence, according to this doctrine, God will also provide the means, circumstances, and situations required to achieve the ultimate end for the elect, or predestined. This is exactly what Guenon was getting at with the notion of “compossible”. The conditions of birth need to be consistent with, and provide the necessary opportunities, to achieve salvation or liberation. A sign of being predestined is to be born into the Medieval Germanic-Roman tradition, or otherwise to have the opportunity to embrace it freely. Nevertheless, some would spurn God's love and reject it.

\paragraph{Antecedent and Consequent Will}
To fully grasp these doctrines, it is necessary to be clear about the difference between the antecedent and consequent will of God. What God wills antecedently may or may not take place. On the other hand, the consequent will is efficacious.

Now the antecedent will of God is for all to achieve salvation, but not all are predestined. The antecedent will is more like the Hermetic notion of Providence. To deny the antecedent will would make God responsible for “wars, concentration camps, and physical and psychical epidemics” as Tomberg points out. This view of God's omnipotence is not uncommon and is based on presumption.

The consequent will involves the alliance of the divine will and the human will. This is not pure passivity, as Fr Garrigou-Lagrange makes clear:

\begin{quotex}
The wills of men are more in God's power than their own.

\end{quotex}
We can't go into all the details right now, but God affects the will through graces and virtues.

\paragraph{Postscript on Incarnation}
Tomberg's understanding of incarnation is in general agreement with Evola's. Tomberg writes (in Letter XX):

\begin{quotex}
The individuality descends consciously and of his own free will to birth, into an environment where he is wanted and awaited.

\end{quotex}
Those are rarer cases however. For most, the incarnation is dominated more by the horizontal than the vertical.

The consequences of Tomberg's insight have been mentioned in other contexts. Some are “born from above”, even if their memory of it is quite dim. The second birth is a remembrance of that. Many, if not most, people today, have a self-understanding of being a product of biological forces and cultural influences. They lack that true will that Tomberg mentions.

\begin{quotex}
The united will constitutes the indestructible and immortal kernel of the body … its active principle, its formative will-energy, survives death.

\end{quotex}
The notion of being “adopted sons [or children] of God” is trivialized today. It leads to a spiritual infantilisation, a pure passivity, the presumption that God loves you just the way you are, without any demands. As we have seen, God does not love everyone to the same degree. His real children are those who have taken on his nature. That's what being a son means: to inherit the nature of the father, or, speaking colloquially, being a “chip off the old block.”

Nevertheless, unlike certain reformation heresies, the Germanic-Roman religion does not accept that some people are condemned at birth. God does not demand the impossible, so it is incorrect to claim that a man may be born with certain proclivities that compel him to deviate from the cosmic law. Even worse is to believe that such compulsions are therefore good and defensible.

\flrightit{Posted on 2015-07-10 by Cologero}
